\chapter{Interaction bidirectionnelle : conjugaison des propagateurs avancé et retardé et conditions de bord de l'absorbeur}
\label{chap:mt-accion-interaccion-absorbedor}

Un chemin devient dynamique lorsqu'il reçoit une fonctionnelle d'action et des
conditions de frontière. L'action attribue un poids à l'histoire complète ; sa
variation sélectionne des transports compatibles avec les extrémités ;
l'interaction couple des degrés de liberté qui partagent une frontière ; la lecture
bidirectionnelle conserve les orientations temporelles directe et inverse du
même antécédent. Dans ce cadre, émission et absorption forment les extrémités d'un
problème aux limites commun.

L'absorbeur occupe une position précise dans l'architecture. HMT apporte le
ruban orienté, la mémoire d'aller et de retour et les involutions discrètes.
La mécanique dimensionnelle ajoute l'opérateur d'évolution, les courants et le
transducteur d'action. Une réalisation physique d'absorbeur relie cette
structure aux propagateurs avancés et retardés et aux conditions de
bord causales. Les identités algébriques du premier tronçon sont
démontrées dans leurs domaines ; le dernier lien conserve le statut
d'interface physique jusqu'à la construction de l'entrelaceur complet.

\section{Fonctionnelle d'action sur des chemins enrichis}

Soit $\gamma=((c_0,\mu_0),(d_1,\ldots,d_L))$ un chemin fini, avec des directions
$d_t$ et des feuillets $\mu_t\in\{\Sigma,\Pi\}$. Nous définissons
\[
 \operatorname{turn}(\gamma)
 =\#\{2\leq t\leq L:d_t\neq d_{t-1}\},
\]
\[
 \operatorname{switch}(\gamma)
 =\#\{2\leq t\leq L:\mu_t\neq\mu_{t-1}\}.
\]

\begin{definicionnarrativa}[title={Action discrète d'interaction minimale}]
\label{def:mt-accion-minima}
Pour $\lambda\geq0$, l'action locale de chemin est
\begin{equation}
 \mathcal A_\lambda(\gamma)
 =\operatorname{turn}(\gamma)
 +\lambda\operatorname{switch}(\gamma)
 =\sum_{t=2}^{L}
 \bigl[1-\delta_{d_t,d_{t-1}}
 +\lambda(1-\delta_{\mu_t,\mu_{t-1}})\bigr].
 \label{eq:mt-accion-minima}
\end{equation}
La première contribution mesure la variation directionnelle et la deuxième mesure le
transport entre feuillets.
\end{definicionnarrativa}

L'action générale incorpore aussi bord, retenue et mémoire :
\begin{equation}
 \mathcal S[\gamma]
 =\sum_{e\in\gamma}\mathcal L_e
 +\mathcal B_{\rm in}(\partial_-\gamma)
 +\mathcal B_{\rm out}(\partial_+\gamma)
 +\mathcal M(\gamma).
 \label{eq:mt-accion-ruta-general}
\end{equation}
Ici, $\mathcal L_e$ est le coût local, $\mathcal B_{\rm in/out}$ fixe les
données admissibles des extrémités et $\mathcal M$ conserve l'information que
la publication visible contracte. La variation s'effectue dans une classe
de chemins ayant les mêmes types de frontière.

\begin{lecturafisica}[title={Principe d'interaction minimale compatible avec la mémoire}]
\label{prin:mt-minima-interaccion}
Parmi les prolongements compatibles avec les données de frontière, la
réalisation physique sélectionnée rend extrémale la fonctionnelle d'action dans le
secteur qui conserve la mémoire exigée par la composition.
\end{lecturafisica}

Le principe fixe la classe du problème variationnel. Une dynamique concrète
ajoute le Hamiltonien ou le lagrangien du secteur et démontre que son extremum
produit les équations du mouvement. La mémoire fonctionne comme contrainte
dynamique : deux chemins ayant la même ombre finale peuvent appartenir à des classes
variationnelles distinctes.

\section{Interaction comme couplage typé}

Soient $\mathfrak p_1$ et $\mathfrak p_2$ deux sections réalisées sur des chemins
$\gamma_1$ et $\gamma_2$. Une interaction à une frontière $B$ consiste en
un morphisme
\begin{equation}
 \mathcal I_B:
 E^{(1)}_B\otimes E^{(2)}_B
 \longrightarrow E^{(12)}_B
 \label{eq:mt-interaccion-frontera}
\end{equation}
et en un terme d'action $\mathcal S_{\rm int}$ compatible avec les
représentations du secteur. Le morphisme conserve les types partagés et
transforme les degrés internes selon les règles de couplage. La
composition complète prend la forme
\[
 (\gamma_2,s_2)\star_{(B,\mathcal I_B)}(\gamma_1,s_1),
\]
où le registre de $B$ est compté une fois et la mémoire des deux côtés est
transportée vers l'état composé.

La covariance de l'interaction exige, pour chaque transformation de symétrie
$g$,
\begin{equation}
 \rho_{12}(g)\,\mathcal I_B
 =\mathcal I_B\bigl(\rho_1(g)\otimes\rho_2(g)\bigr).
 \label{eq:mt-covariancia-interaccion}
\end{equation}
Cette identité décide si le couplage appartient à la représentation
physique déclarée. Dans les secteurs électromagnétique et de couleur, la même
structure s'exprime au moyen d'une connexion et de son transport parallèle.

\section{Conjugaison temporelle du propagateur}

Soit $\mathcal H_{\mathbb R}$ la réalification de l'espace quantique, munie
d'une structure complexe $J_E$ avec $J_E^2=-I$. Soit $H=H^*$ un opérateur
autoadjoint à domaine commun invariant et soit $C$ une involution orthogonale
qui satisfait
\begin{equation}
 HJ_E=J_EH,\qquad
 CJ_EC^{-1}=-J_E,
 \qquad CHC^{-1}=H.
 \label{eq:mt-hipotesis-reversion-propagador}
\end{equation}

\begin{resultado}[title={Théorème : Réversion bidirectionnelle du groupe d'évolution}]
\label{thm:mt-reversion-propagador}
L'opérateur
\begin{equation}
 U(t)=\exp\!\left(-\frac{J_EHt}{\hbar_{\rm int}}\right)
 \label{eq:mt-grupo-evolucion}
\end{equation}
forme un groupe unitaire à un paramètre et satisfait
\begin{equation}
 \boxed{CU(t)C^{-1}=U(-t)}.
 \label{eq:mt-reversion-unitaria}
\end{equation}
\end{resultado}

\paragraph{Démonstration de la réversion bidirectionnelle du groupe unitaire \(U(t)\) : \(CU(t)C^{-1}=U(-t)\)}
Le générateur $A=-J_EH/\hbar_{\rm int}$ est antiadjoint. Le théorème de Stone
produit le groupe unitaire $U(t)=e^{tA}$. Les hypothèses de
\eqref{eq:mt-hipotesis-reversion-propagador} donnent $CAC^{-1}=-A$ ; le calcul
fonctionnel fournit $Ce^{tA}C^{-1}=e^{-tA}$.


L'involution $C_M=-\operatorname{rev}$ des mots de chemin apporte la
donnée généalogique d'orientation. Lorsque le foncteur physique envoie $C_M$ sur
l'involution $C$ du théorème, les sections particule et antiparticule sont
représentées par des orientations temporelles conjuguées du même groupe
unitaire. L'identité est démontrée avec la structure de départ
explicite : autoadjonction, domaine, structure complexe et involution.

\section{Propagateurs avancés et retardés}

Soit $L$ l'opérateur hyperbolique d'une réalisation de champs et soient
$G_{\rm ret}$ et $G_{\rm adv}$ ses inverses fondamentaux à supports
causaux futur et passé, respectivement :
\begin{equation}
 LG_{\rm ret}=LG_{\rm adv}=I,
 \qquad
 \operatorname{supp}(G_{\rm ret}f)\subset J^+(\operatorname{supp}f),
 \qquad
 \operatorname{supp}(G_{\rm adv}f)\subset J^-(\operatorname{supp}f).
 \label{eq:mt-green-avanzado-retardado}
\end{equation}
La différence causale
$E=G_{\rm ret}-G_{\rm adv}$ transporte des données entre solutions et la moyenne
symétrique
\[
 G_{\rm sym}=\frac12(G_{\rm ret}+G_{\rm adv})
\]
réalise un noyau temporellement symétrique.

Une interface HMT--MD vers cette théorie est une famille d'applications
\begin{equation}
 \mathfrak I_{\rm prop}:
 (\gamma,\gamma^\vee,\mathcal B_\gamma,U,C_M)
 \longrightarrow
 (G_{\rm ret},G_{\rm adv},E,\mathcal H_{\rm field})
 \label{eq:mt-interfaz-green}
\end{equation}
qui entrelace l'involution de chemin avec l'échange
$G_{\rm ret}\leftrightarrow G_{\rm adv}$, préserve le produit intérieur ou la
forme symplectique pertinente et envoie les actions des deux côtés sur la même
forme fonctionnelle. Cette application constitue la condition mathématique qui transforme
l'orientation bidirectionnelle interne en propagation avancée--retardée de
champs.

\section{Réalisation HMT–MD de l'absorbeur par inversion des chemins et conjugaison particule–antiparticule}

\begin{definicionnarrativa}[title={Système d'émission--absorption bidirectionnel}]
\label{def:mt-absorbedor}
Une réalisation physique d'absorbeur est un problème aux limites formé par
un courant émetteur $J_{\rm em}$, une réponse absorbante $J_{\rm abs}$,
les propagateurs $G_{\rm ret}$ et $G_{\rm adv}$ et des conditions de bord qui
ferment conjointement les deux orientations du transport. Son action
effective temporellement symétrique possède la forme
\begin{equation}
 \mathcal S_{\rm abs}[J]
 =\mathcal S_{\rm matter}[J]
 +\frac12\langle J,G_{\rm sym}J\rangle,
 \qquad J=J_{\rm em}+J_{\rm abs},
 \label{eq:mt-accion-absorbedor}
\end{equation}
et ses conditions aux limites déterminent la partie radiative observée.
\end{definicionnarrativa}

La définition exprime une théorie d'action à distance ou de réponse aux
limites en langage opératoire. Dans HMT--MD, le ruban apporte la paire
$(\gamma,\gamma^\vee)$ ; le résultat de réversion bidirectionnelle apporte
la réversion exacte du groupe unitaire ; le foncteur
$\mathfrak I_{\rm prop}$ doit apporter causalité, produit intérieur et
conditions de bord. L'architecture de l'absorbeur prolonge la structure
bidirectionnelle déjà construite. Sa réalisation reste une interface physique,
et la construction de l'entrelaceur causal complet constitue un problème
mathématique explicitement délimité.

La clôture physique exige quatre identités vérifiables :
\begin{enumerate}
 \item commutativité entre conjugaison de chemin et échange des deux
 propagateurs ;
 \item égalité des actions transportées par
 $\mathfrak I_{\rm prop}$ ;
 \item conservation de la forme symplectique ou du produit intérieur ;
 \item compatibilité des conditions d'émission et d'absorption sous
 raffinement du chemin.
\end{enumerate}
Un échec de l'une quelconque de ces identités réfute la réalisation particulière
et laisse intacte la réversion opératoire déjà démontrée dans
\eqref{eq:mt-reversion-unitaria}.

\section{Les actions séparées de \texorpdfstring{$C$}{C},
\texorpdfstring{$P$}{P} et \texorpdfstring{$T$}{T}}

Sur un chemin complet, nous définissons trois involutions typées :
\begin{enumerate}
 \item $C$ échange les feuillets $\Sigma\leftrightarrow\Pi$ et conjugue les
 caractères orientés ;
 \item $P$ réfléchit la carte spatiale, par exemple
 $(i,j)\mapsto(i,-j)\pmod9$, avec
 $E\leftrightarrow W$ ;
 \item $T$ inverse l'ordre du registre, remplace chaque direction par son
 opposée et permute les extrémités du chemin.
\end{enumerate}
Chaque opération agit sur une composante différente de l'état. Leur composition
$CPT$ agit sur le feuillet, l'espace et l'orientation temporelle à la fois.

\begin{resultado}[title={Théorème : Invariance discrète individuelle et conjointe}]
\label{thm:mt-cpt-rutas}
Pour tout chemin $\gamma$ et tout $\lambda\geq0$,
\begin{equation}
 \mathcal A_\lambda(\gamma)
 =\mathcal A_\lambda(C\gamma)
 =\mathcal A_\lambda(P\gamma)
 =\mathcal A_\lambda(T\gamma)
 =\mathcal A_\lambda(CPT\gamma).
 \label{eq:mt-invariancia-cpt-rutas}
\end{equation}
\end{resultado}

\paragraph{Compatibilité des actions \(C\), \(P\) et \(T\) avec la réversion unitaire \(U(t)\mapsto U(-t)\)}
$P$ permute l'alphabet des directions et conserve le motif des égalités
entre pas consécutifs. $C$ permute globalement les feuillets et conserve le
motif des changements de feuillet. $T$ inverse l'ordre des deux suites et
conserve le nombre de transitions adjacentes. Chaque involution préserve
séparément les deux termes de
\eqref{eq:mt-accion-minima} ; la composition hérite de l'invariance.


Le théorème appartient au registre discret et à son action locale. La matrice CKM
réalisée dans la base de saveur possède un invariant de Jarlskog différent de
zéro et donc une violation CP dans cette représentation. Les deux résultats sont
compatibles : un lecteur orienté peut distinguer des états reliés par
$CP$ tandis que la fonctionnelle locale
$\mathcal A_\lambda$ conserve chaque involution. La
réalisation du théorème CPT relativiste des champs ajoute localité,
covariance de Lorentz, structure spectrale et représentation de champs. Dans
ce domaine élargi, la CPT conjointe constitue la symétrie à préserver. Dans le
domaine interne, la clôture CKM et l'expansion temporelle restent compatibles
avec l'invariance CPT démontrée pour l'action discrète.

Trois strates sont ainsi séparées. L'invariance de
\eqref{eq:mt-invariancia-cpt-rutas} est un théorème interne exact. L'identité
\eqref{eq:mt-reversion-unitaria} est démontrée avec une structure de départ
explicite. La théorie physique de l'absorbeur et le théorème CPT des champs exigent
les entrelaceurs propres à leurs codomaines. Cette séparation maintient
l'orientation temporelle inverse dans la définition de particule et accorde
à chaque affirmation la force qui lui correspond.
